\section{PHÂN TÍCH HỆ THỐNG VÀ THUYẾT MINH GIẢI PHÁP KỸ THUẬT
}
\subsection{Phân tích hệ thống}
\subsubsection{Mô tả quy trình nghiệp vụ}
\begin{figure}[htpb]
    \centering
    % Resizebox giúp sơ đồ tự động co giãn vừa vặn với chiều ngang trang giấy
    \resizebox{\textwidth}{!}{
    \begin{tikzpicture}[
        % Định nghĩa style cho các khối hình
        box/.style={rectangle, draw, thick, minimum width=2.8cm, minimum height=1.2cm, text width=2.5cm, align=center, fill=white, font=\small},
        dia/.style={diamond, draw, thick, minimum width=3cm, minimum height=2.5cm, text width=1.8cm, align=center, fill=white, font=\small},
        line/.style={draw, thick, -Stealth},
        labeltext/.style={font=\scriptsize, align=center}
    ]

        % 1. VẼ KHUNG SWIMLANE (LÀN ĐƯỜNG)
        % Khung ngoài cùng
        \draw[thick] (-9, 0.5) rectangle (9, 12);
        % Đường kẻ ngang tiêu đề
        \draw[thick] (-9, 11) -- (9, 11);
        % Đường kẻ dọc chia 2 làn
        \draw[thick] (0, 0.5) -- (0, 12);
        
        % Tiêu đề của 2 làn
        \node[font=\bfseries] at (-4.5, 11.5) {Hệ thống};
        \node[font=\bfseries] at (4.5, 11.5) {Cán bộ quản lý đô thị};

        % 2. ĐẶT CÁC KHỐI NODE VÀO VỊ TRÍ
        % Làn phải (Cán bộ)
        \node[dia] (loai_tac_vu) at (2, 9) {Loại tác vụ};
        \node[box] (nhap_moi) at (7, 9) {Nhập mới};
        \node[box] (cap_nhat) at (2, 5) {Cập nhật};
        \node[box] (nhap_thong_tin) at (7, 3.85) {Nhập thông tin mô tả};
        \node[box] (gui_yeu_cau) at (7, 2) {Gửi yêu cầu phê duyệt};

        % Làn trái (Hệ thống)
        \node[dia] (du_lieu_hop_le) at (-3, 9) {Dữ liệu hợp lệ};
        \node[box] (luu_tru) at (-7, 9) {Lưu trữ vào CSDL};
        \node[box] (yeu_cau_chinh_sua) at (-3, 3.5) {Yêu cầu chỉnh sửa nội dung};

        % 3. VẼ CÁC ĐƯỜNG MŨI TÊN LIÊN KẾT
        % Các mũi tên trong làn phải
        \draw[line] (loai_tac_vu.east) -- node[above, labeltext] {Nhập thông\\tin tài sản} (nhap_moi.west);
        \draw[line] (loai_tac_vu.south) -- node[right, labeltext] {Đo đạc thực\\địa/Cập nhật} (cap_nhat.north);
        \draw[line] (nhap_moi.south) -- (nhap_thong_tin.north);
        \draw[line] (cap_nhat.east) -- (nhap_thong_tin.west);
        \draw[line] (nhap_thong_tin.south) -- (gui_yeu_cau.north);

        % Mũi tên từ làn phải sang làn trái (Gửi yêu cầu -> Dữ liệu hợp lệ)
        % (Đi sang trái tới tọa độ X = -0.5, sau đó đi lên và rẽ trái vào khối)
        \draw[line] (gui_yeu_cau.west) -- (-0.5, 2) |- (du_lieu_hop_le.east);

        % Các mũi tên trong làn trái
        \draw[line] (du_lieu_hop_le.west) -- node[above, labeltext] {Đúng} (luu_tru.east);
        \draw[line] (du_lieu_hop_le.south) -- node[left, labeltext] {Từ chối} (yeu_cau_chinh_sua.north);

        % Mũi tên từ làn trái trả về làn phải (Yêu cầu chỉnh sửa -> Nhập thông tin)
        % (Đi sang phải tới tọa độ X = 1.2, sau đó đi lên và rẽ phải vào khối)
        % [yshift=-10pt] giúp mũi tên đi vào dưới đường của khối Cập nhật để không bị đè nét
        \draw[line] (yeu_cau_chinh_sua.east) -- (1.2, 3.5) |- ([yshift=-10pt]nhap_thong_tin.west);

    \end{tikzpicture}
    }
    \caption{Sơ đồ quy trình cập nhật thông tin tài sản}
    \label{fig:quy_trinh_cap_nhat}
\end{figure}
\textbf{Mô tả quy trình:}
\begin{enumerate}
    \item \textbf{Bước 1:} Cán bộ quản lý đô thị thực hiện khởi tạo dữ liệu dựa trên loại tác vụ:
    \begin{itemize}
        \item Nếu là \textit{Nhập thông tin tài sản}: Cán bộ chọn tác vụ \textbf{Nhập mới}.
        \item Nếu là \textit{Đo đạc thực địa/Cập nhật}: Cán bộ chọn tác vụ \textbf{Cập nhật}.
    \end{itemize}
    
    \item \textbf{Bước 2:} Sau đó, cán bộ tiến hành \textbf{Nhập thông tin mô tả} chi tiết cho dữ liệu/tài sản đó.
    
    \item \textbf{Bước 3:} Cán bộ \textbf{Gửi yêu cầu phê duyệt} để hệ thống kiểm tra dữ liệu có hợp lệ hay không.
    
    \item \textbf{Bước 4:} Hệ thống kiểm tra tính hợp lệ của dữ liệu và đưa ra quyết định theo 2 hướng:
    \begin{itemize}
        \item \textbf{Trường hợp Đúng:} Dữ liệu sẽ được hệ thống tự động \textbf{Lưu trữ vào \gls{csdl}}.
        \item \textbf{Trường hợp Sai (Từ chối):} Hệ thống gửi \textbf{Yêu cầu chỉnh sửa nội dung}. Khi đó, quy trình sẽ quay trở lại \textbf{Bước 2} để Cán bộ quản lý đô thị kiểm tra và cập nhật lại thông tin mô tả.
    \end{itemize}
\end{enumerate}
\subsubsection{ Quy trình nghiệp vụ Thống kê và phân tích dữ liệu}
\begin{figure}[htbp]
  \centering
  % Tự động thu phóng sơ đồ vừa với chiều rộng trang giấy
  \resizebox{\textwidth}{!}{
  \begin{tikzpicture}[
    % Định nghĩa khoảng cách và kiểu dáng các khối
    block/.style={rectangle, draw, thick, minimum width=3.5cm, minimum height=1.2cm, text width=3.2cm, align=center, fill=white, font=\small},
    decision/.style={diamond, draw, thick, minimum width=3cm, minimum height=2.5cm, text width=2cm, align=center, fill=white, font=\small},
    arrow/.style={draw, thick, -Stealth},
    ]

    % 1. VẼ KHUNG SWIMLANE
    % Khung bao quanh toàn bộ sơ đồ
    \draw [thick] (-7.5, 0) rectangle (7.5, 10.5);
    % Đường kẻ dọc chia 2 làn
    \draw [thick] (0, 0) -- (0, 10.5);
    % Đường kẻ ngang chia phần Tiêu đề
    \draw [thick] (-7.5, 9.5) -- (7.5, 9.5);

    % Tiêu đề của 2 làn
    \node[font=\bfseries] at (-3.75, 10) {Cán bộ quản lý};
    \node[font=\bfseries] at (3.75, 10) {Cơ quan nhà nước};

    % 2. ĐẶT CÁC KHỐI NODE
    % Làn trái (Cán bộ quản lý)
    \node [block] (tong_hop) at (-3.75, 8) {Tổng hợp thông tin};
    \node [block] (thong_ke) at (-3.75, 6) {Thống kê};
    \node [block] (phan_tich) at (-3.75, 4) {Phân tích};
    \node [block] (viet_bao_cao) at (-3.75, 1.5) {Viết báo cáo định hướng phát triển đô thị};

    % Làn phải (Cơ quan nhà nước)
    \node [decision] (phe_duyet) at (3.75, 6) {Phê duyệt};
    \node [block] (ban_hanh) at (3.75, 3) {Ban hành chủ\\trương};

    % 3. VẼ CÁC ĐƯỜNG MŨI TÊN
    % Các mũi tên nối thẳng ở làn trái
    \draw [arrow] (tong_hop.south) -- (thong_ke.north);
    \draw [arrow] (thong_ke.south) -- (phan_tich.north);
    \draw [arrow] (phan_tich.south) -- (viet_bao_cao.north);

    % Mũi tên từ "Viết báo cáo..." sang "Phê duyệt"
    % [yshift=10pt]: Dịch điểm xuất phát lên trên một chút
    % Đi sang phải đến x=-1.2, đi thẳng lên y=8.8, rồi rẽ vuông góc vào đỉnh khối Phê duyệt
    \draw [arrow] ([yshift=10pt]viet_bao_cao.east) -- (-1.2, 1.9) -- (-1.2, 8.8) -| (phe_duyet.north);

    % Mũi tên từ "Phê duyệt" trả về "Viết báo cáo..." (Từ chối)
    % Đi sang trái đến x=-0.5, rẽ vuông góc đi xuống vị trí thấp của khối Viết báo cáo [yshift=-10pt]
\draw [arrow] (phe_duyet.west) -- node[above, pos=0.5] {Từ chối} (-0.5, 6) |- ([yshift=-10pt]viet_bao_cao.east);
    % Mũi tên từ "Phê duyệt" đi xuống "Ban hành" (Đồng ý)
    \draw [arrow] (phe_duyet.south) -- node[right] {Đồng ý} (ban_hanh.north);

  \end{tikzpicture}
  }
  \caption{Quy trình nghiệp vụ Thống kê và phân tích dữ liệu}
  \label{fig:quy_trinh_thong_ke}
\end{figure}
\textbf{Mô tả quy trình:}
\begin{enumerate}
    \item \textbf{Bước 1: Cán bộ quản lý thực hiện chuẩn bị dữ liệu.} Cán bộ quản lý tiến hành quy trình xử lý thông tin qua các giai đoạn:
    \begin{itemize}
        \item \textbf{Tổng hợp thông tin:} Tập hợp các dữ liệu cần thiết từ nhiều nguồn.
        \item \textbf{Thống kê:} Hệ thống hóa số liệu dưới dạng các bảng biểu hoặc chỉ số.
        \item \textbf{Phân tích:} Đánh giá dữ liệu để tìm ra các xu hướng hoặc vấn đề cần giải quyết.
    \end{itemize}
    
    \item \textbf{Bước 2: Lập báo cáo định hướng.} Sau khi có kết quả phân tích, cán bộ thực hiện \textbf{Viết báo cáo định hướng phát triển đô thị} và trình lên cấp có thẩm quyền.
    
    \item \textbf{Bước 3: Cơ quan nhà nước thực hiện phê duyệt.} Cơ quan nhà nước xem xét báo cáo và đưa ra quyết định tại bước \textbf{Phê duyệt}:
    \begin{itemize}
        \item \textbf{Trường hợp Từ chối:} Báo cáo được gửi trả lại cho Cán bộ quản lý để yêu cầu chỉnh sửa, bổ sung hoặc viết lại định hướng.
        \item \textbf{Trường hợp Đồng ý:} Cơ quan chức năng tiến hành \textbf{Ban hành chủ trương} chính thức để triển khai các bước tiếp theo trong thực tế.
    \end{itemize}
\end{enumerate}
\subsection{Đề xuất Quy trình tin học hóa}
\subsubsection{Quy trình tin học hóa Quản lý tài sản đô thị}
\begin{figure}[htbp]
  \centering
  \resizebox{\textwidth}{!}{
  \begin{tikzpicture}[
    % Định nghĩa kiểu dáng khối chữ nhật và hình thoi
    block/.style={rectangle, draw, thick, minimum width=3.8cm, minimum height=1.2cm, text width=3.5cm, align=center, fill=white, font=\small},
    decision/.style={diamond, draw, thick, minimum width=3cm, minimum height=2.5cm, text width=2cm, align=center, fill=white, font=\small},
    arrow/.style={draw, thick, -Stealth},
    ]

    % 1. VẼ KHUNG SWIMLANE
    % Khung bao quanh
    \draw [thick] (-9.5, -1) rectangle (9.5, 10);
    % Đường kẻ dọc chia 2 làn
    \draw [thick] (0, -1) -- (0, 10);
    % Đường kẻ ngang chia tiêu đề
    \draw [thick] (-9.5, 9) -- (9.5, 9);

    % Tiêu đề các làn
    \node[font=\bfseries] at (-4.75, 9.5) {Hệ thống};
    \node[font=\bfseries] at (4.75, 9.5) {Cán bộ quản lý đô thị};

    % 2. ĐẶT CÁC NODE KHỐI
    % Làn trái (Hệ thống)
    \node [block] (ai_nhan_dien) at (-6, 7.5) {\gls{ai} nhận diện các cơ sở hạ tầng có trên bản đồ \gls{gis}};
    \node [decision] (hop_le) at (-7, 3) {Dữ liệu hợp lệ};
    \node [block] (yeu_cau) at (-2.1, 3) {Yêu cầu chỉnh sửa nội dung};
    \node [block] (luu_tru) at (-7, 0) {Lưu trữ vào \gls{csdl}};

    % Làn phải (Cán bộ)
    \node [block] (kiem_tra) at (5, 7.5) {Kiểm tra và chỉnh sửa thông tin};
    \node [block] (gui_yeu_cau) at (5, 5.5) {Gửi yêu cầu phê duyệt};

    % 3. VẼ CÁC ĐƯỜNG MŨI TÊN
    % AI nhận diện -> Kiểm tra (Mũi tên đi ngang vào nửa trên khối Kiểm tra)
    \draw [arrow] ([yshift=10pt]ai_nhan_dien.east) -- ([yshift=10pt]kiem_tra.west);

    % Kiểm tra -> Gửi yêu cầu
    \draw [arrow] (kiem_tra.south) -- (gui_yeu_cau.north);

    % Gửi yêu cầu -> Dữ liệu hợp lệ
    % Cú pháp -| giúp mũi tên đi ngang sang trái rồi bẻ góc vuông đi xuống
    \draw [arrow] (gui_yeu_cau.west) -| (hop_le.north);

    % Dữ liệu hợp lệ -> Lưu trữ (Nhánh Đúng)
    \draw [arrow] (hop_le.south) -- node[right] {Đúng} (luu_tru.north);

    % Dữ liệu hợp lệ -> Yêu cầu chỉnh sửa (Nhánh Sai)
    \draw [arrow] (hop_le.east) -- node[above] {Sai} (yeu_cau.west);

    % Yêu cầu chỉnh sửa -> Kiểm tra (Mũi tên trả về)
    % Cú pháp |- giúp mũi tên đi dọc lên trên rồi bẻ góc vuông sang phải vào nửa dưới khối Kiểm tra
    \draw [arrow] (yeu_cau.north) |- ([yshift=-10pt]kiem_tra.west);

  \end{tikzpicture}
  }
  \caption{Quy trình tin học hóa Quản lý tài sản đô thị}
  \label{fig:quy_trinh_tin_hoc_hoa}
\end{figure}
\textbf{Mô tả quy trình:}
\begin{enumerate}
    \item \textbf{Bước 1: Hệ thống tự động nhận diện dữ liệu.} \\
    Hệ thống sử dụng công nghệ ``\gls{ai} nhận diện các cơ sở hạ tầng có trên bản đồ \gls{gis}''. Sau khi nhận diện, hệ thống chuyển kết quả cho bộ phận chuyên môn.
    
    \item \textbf{Bước 2: Cán bộ quản lý đô thị xử lý thông tin.} \\
    Cán bộ thực hiện \textbf{Kiểm tra và chỉnh sửa thông tin} để xác minh độ chính xác của các đối tượng mà \gls{ai} đã nhận diện. Tiếp theo, cán bộ tiến hành \textbf{Gửi yêu cầu phê duyệt} lên hệ thống.
    
    \item \textbf{Bước 3: Hệ thống kiểm tra tính hợp lệ của dữ liệu và đưa ra quyết định theo 2 hướng:}
    \begin{itemize}
        \item \textbf{Trường hợp Đúng:} Dữ liệu sẽ được hệ thống tự động \textbf{Lưu trữ vào \gls{csdl}}.
        \item \textbf{Trường hợp Sai (Từ chối):} Hệ thống gửi \textbf{Yêu cầu chỉnh sửa nội dung}. Khi đó, quy trình sẽ quay trở lại Bước 2 để Cán bộ quản lý đô thị kiểm tra và cập nhật lại thông tin mô tả.
    \end{itemize}
\end{enumerate}
\subsubsection{Quy trình tin học hóa Người dân phản ánh sự cố}
\begin{figure}[htbp]
  \centering
  % Tự động thu phóng vừa với chiều ngang trang giấy
  \resizebox{\textwidth}{!}{
  \begin{tikzpicture}[
    % Hạ cỡ chữ xuống \footnotesize, tăng chiều rộng khối lên 2.8cm
    block/.style={rectangle, draw, thick, minimum width=2.8cm, minimum height=1.2cm, text width=2.6cm, align=center, fill=white, font=\footnotesize},
    decision/.style={diamond, draw, thick, minimum width=2.5cm, minimum height=2.5cm, text width=1.8cm, align=center, fill=white, font=\footnotesize},
    arrow/.style={draw, thick, -Stealth},
    ]

    % 1. VẼ KHUNG SWIMLANE (LÀN ĐƯỜNG)
    % Khung bao ngoài cùng (Mở rộng bao trọn toàn bộ các khối)
    \draw [thick] (-5.5, -2) rectangle (9, 10.5);
    
    % Đường kẻ dọc chia 3 làn (Mở rộng khoảng cách các làn)
    \draw [thick] (-1.5, -2) -- (-1.5, 10.5);
    \draw [thick] (4, -2) -- (4, 10.5);
    
    % Đường kẻ ngang chia phần Tiêu đề
    \draw [thick] (-5.5, 9.5) -- (9, 9.5);

    % Tiêu đề của 3 làn
    \node[font=\bfseries] at (-3.5, 10) {Người dân};
    \node[font=\bfseries] at (1.25, 10) {Hệ thống};
    \node[font=\bfseries] at (6.5, 10) {Đơn vị xử lý};

    % 2. ĐẶT CÁC KHỐI NODE THEO TỌA ĐỘ
    % Làn 1: Người dân (Tâm X = -3.5)
    \node [block] (phat_hien) at (-3.5, 8) {Phát hiện sự cố};
    \node [block] (phan_anh) at (-3.5, 6) {Phản ánh sự cố};

    % Làn 2: Hệ thống (Tâm X = 1.25)
    \node [block] (tiep_nhan) at (1.25, 8) {Tự động tiếp nhận sự cố};
    \node [block] (phan_loai) at (1.25, 6) {Phân loại mức độ ưu tiên};
    \node [block] (cap_nhat_td) at (1.25, 3) {Cập nhật tiến độ};
    \node [block] (luu_tru) at (1.25, -0.5) {Lưu trữ vào \gls{csdl}};

    % Làn 3: Đơn vị xử lý (Tâm X = 6.5)
    \node [block] (xac_thuc) at (6.5, 8) {Xác thực thông tin};
    \node [block] (thuc_hien) at (6.5, 6) {Thực hiện xử lý sự cố};
    \node [decision] (hoan_thanh) at (6.5, 3) {Hoàn thành};
    \node [block] (cap_nhat_tt) at (6.5, -0.5) {Cập nhật thông tin cơ sở hạ tầng};


    % 3. VẼ CÁC ĐƯỜNG MŨI TÊN
    % --- Nội bộ làn 1 ---
    \draw [arrow] (phat_hien) -- (phan_anh);

    % --- Làn 1 sang Làn 2 ---
    % Đi ngang sang phải (thêm 1cm) rồi gấp khúc đi thẳng lên khối Tiếp nhận
    \draw [arrow] (phan_anh.east) -- ++(1,0) |- (tiep_nhan.west);

    % --- Nội bộ làn 2 ---
    \draw [arrow] (tiep_nhan) -- (phan_loai);

    % --- Làn 2 sang Làn 3 ---
    \draw [arrow] (phan_loai.east) -- ++(1,0) |- (xac_thuc.west);

    % --- Nội bộ làn 3 ---
    \draw [arrow] (xac_thuc) -- (thuc_hien);
    \draw [arrow] (thuc_hien) -- (hoan_thanh);

    % --- Trả kết quả từ Quyết định ---
    % Trường hợp Sai: Từ Làn 3 chéo ngang về Làn 2
    \draw [arrow] (hoan_thanh.west) -- node[above,pos=0.6] {Sai} (cap_nhat_td.east);
    
    % Vòng lặp từ Cập nhật tiến độ (Làn 2) quay lại Thực hiện xử lý (Làn 3)
    % Cú pháp này đi lên trên 1 đoạn, rẽ ngang sang phải, rồi chui vào ĐÁY của khối "Thực hiện xử lý" để không đè lên mũi tên Phân loại
    \draw [arrow] (cap_nhat_td.north) -- ++(0, 1.2) -| ([xshift=-15pt]thuc_hien.south);

    % Trường hợp Đúng: Đi thẳng xuống
    \draw [arrow] (hoan_thanh.south) -- node[right] {Đúng} (cap_nhat_tt.north);

    % --- Từ Làn 3 về lại Làn 2 (Lưu trữ) ---
    \draw [arrow] (cap_nhat_tt.west) -- (luu_tru.east);

  \end{tikzpicture}
  }
  \caption{Quy trình tin học hóa Người dân phản ánh sự cố}
  \label{fig:quy_trinh_phan_anh}
\end{figure}
\textbf{Mô tả quy trình:}
\textbf{Mô tả quy trình:}
\begin{itemize}
    \item \textbf{Bước 1: Người dân thực hiện gửi phản ánh.} \\
    Khi phát hiện sự cố, người dân thực hiện gửi thông tin qua bước \textit{Phản ánh sự cố} đến hệ thống quản lý.
    
    \item \textbf{Bước 2: Hệ thống tiếp nhận và phân loại tự động.} \\
    Hệ thống thực hiện \textit{Tự động tiếp nhận sự cố} từ phản ánh của người dân. Sau đó, hệ thống tiến hành \textit{Phân loại mức độ ưu tiên} để chuyển thông tin đến đơn vị có thẩm quyền xử lý tương ứng.
    
    \item \textbf{Bước 3: Đơn vị xử lý thực hiện nghiệp vụ.} \\
    Đơn vị xử lý tiến hành \textit{Xác thực thông tin} để kiểm tra tính chính xác của sự cố. Sau khi xác thực, đơn vị thực hiện bước \textit{Thực hiện xử lý sự cố}.
    
    \item \textbf{Bước 4: Kiểm tra kết quả và lưu trữ dữ liệu.} \\
    Trạng thái công việc được đánh giá tại bước \textit{Hoàn thành}:
    \begin{itemize}
        \item \textbf{Trường hợp chưa hoàn thành (Sai):} Đơn vị thực hiện \textit{Cập nhật tiến độ} lên hệ thống và tiếp tục quay lại quy trình xử lý.
        \item \textbf{Trường hợp đã hoàn thành (Đúng):} Đơn vị tiến hành \textit{Cập nhật thông tin cơ sở hạ tầng}. Cuối cùng, thông tin được chuyển về hệ thống để thực hiện \textit{Lưu trữ vào \gls{csdl}}, hoàn tất quy trình.
    \end{itemize}
\end{itemize}
\subsection{Quy trình tin học hóa Thống kê và phân tích dữ liệu}
\begin{figure}[htbp]
  \centering
  \resizebox{\textwidth}{!}{
  \begin{tikzpicture}[
    % Định nghĩa kiểu dáng khối. Dùng \footnotesize để chữ gọn gàng trong hộp
    block/.style={rectangle, draw, thick, minimum width=3.2cm, minimum height=1.2cm, text width=3cm, align=center, fill=white, font=\footnotesize},
    decision/.style={diamond, draw, thick, minimum width=2.8cm, minimum height=2.8cm, text width=1.8cm, align=center, fill=white, font=\footnotesize},
    arrow/.style={draw, thick, -Stealth}
    ]

    % 1. VẼ KHUNG SWIMLANE RỘNG RÃI
    % Tổng chiều rộng từ -6.75 đến 6.75 (chia làm 3 cột, mỗi cột rộng 4.5cm)
    \draw [thick] (-6.75, 0) rectangle (7.25, 9.5);
    \draw [thick] (-2.25, 0) -- (-2.25, 9.5); % Vách 1
    \draw [thick] (2.25, 0) -- (2.25, 9.5);   % Vách 2
    \draw [thick] (-6.75, 8.5) -- (7.25, 8.5); % Vách ngang tiêu đề

    % Tiêu đề các làn (Căn giữa mỗi cột)
    \node[font=\bfseries] at (-4.5, 9) {Hệ thống};
    \node[font=\bfseries] at (0, 9) {Cán bộ quản lý};
    \node[font=\bfseries] at (4.5, 9) {Cơ quan nhà nước};

    % 2. ĐẶT CÁC NODE VÀO ĐÚNG TÂM TỪNG CỘT
    % Làn 1: Hệ thống (Tâm X = -4.5)
    \node [block] (thong_ke) at (-4.5, 7) {Thống kê};

    % Làn 2: Cán bộ quản lý (Tâm X = 0)
    \node [block] (truy_van) at (0, 7) {Truy vấn dữ liệu};
    \node [block] (phan_tich) at (0, 4.5) {Phân tích};
    \node [block] (viet_bao_cao) at (0, 1.5) {Viết báo cáo định hướng phát triển đô thị};

    % Làn 3: Cơ quan nhà nước (Tâm X = 4.5)
    % Đẩy khối Phê duyệt lên cao một chút (Y=6) để lấy không gian cho mũi tên chạy
    \node [decision] (phe_duyet) at (5.5, 6) {Phê duyệt}; 
    \node [block] (ban_hanh) at (5.5, 1.5) {Ban hành chủ trương};

    % 3. VẼ ĐƯỜNG MŨI TÊN
    % Làn 1 -> Làn 2
    \draw [arrow] (thong_ke) -- (truy_van);
    
    % Nội bộ Làn 2
    \draw [arrow] (truy_van) -- (phan_tich);
    \draw [arrow] (phan_tich) -- (viet_bao_cao);

    % Làn 2 -> Làn 3 (Viết báo cáo -> Phê duyệt)
    % Điểm xuất phát nhích lên trên 10pt. Kéo ngang sang phải 1.5cm rồi mới gấp khúc đi thẳng lên đỉnh Phê duyệt
    \draw [arrow] ([yshift=10pt]viet_bao_cao.east) -- ++(0.4,0) |- (phe_duyet.north);

    % Làn 3 -> Làn 2 (Từ chối)
    % Xuất phát từ cạnh trái Phê duyệt. Kéo ngang sang trái 1.5cm rồi gấp khúc đi thẳng xuống vị trí thấp của Viết báo cáo (-10pt)
    \draw [arrow] (phe_duyet.west) -- node[above] {Từ chối} ++(-1.5,0) |- ([yshift=-10pt]viet_bao_cao.east);

    % Nội bộ Làn 3 (Đồng ý)
    \draw [arrow] (phe_duyet.south) -- node[right] {Đồng ý} (ban_hanh.north);

  \end{tikzpicture}
  }
  \caption{Quy trình tin học hóa Thống kê và phân tích dữ liệu}
  \label{fig:quy_trinh_thong_ke_data}
\end{figure}
\textbf{Mô tả quy trình:}
\begin{enumerate}
    \item \textbf{Bước 1: Hệ thống cung cấp dữ liệu nền tảng.} \\
    Hệ thống thực hiện bước \textbf{Thống kê} các số liệu thô và chuyển thông tin cho cán bộ chuyên môn.
    
    \item \textbf{Bước 2: Cán bộ quản lý xử lý và lập báo cáo.} \\
    Cán bộ quản lý tiếp nhận dữ liệu và thực hiện các nghiệp vụ sau:
    \begin{itemize}
        \item \textbf{Truy vấn dữ liệu:} Kiểm tra và trích xuất các thông tin cụ thể cần thiết.
        \item \textbf{Phân tích:} Đánh giá số liệu để tìm ra các định hướng phát triển.
        \item \textbf{Viết báo cáo định hướng phát triển đô thị:} Hoàn thiện văn bản báo cáo và trình lên cơ quan cấp trên.
    \end{itemize}
    
    \item \textbf{Bước 3: Cơ quan nhà nước thực hiện phê duyệt.} \\
    Cơ quan có thẩm quyền tiến hành xem xét báo cáo tại bước \textbf{Phê duyệt}:
    \begin{itemize}
        \item \textbf{Trường hợp Từ chối:} Báo cáo được gửi trả lại để Cán bộ quản lý chỉnh sửa và cập nhật lại nội dung định hướng.
        \item \textbf{Trường hợp Đồng ý:} Cơ quan nhà nước tiến hành \textbf{Ban hành chủ trương} chính thức để làm căn cứ thực hiện các kế hoạch phát triển đô thị.
    \end{itemize}
\end{enumerate}

\subsection{Xác định khối lượng công việc}
Khối lượng công việc toàn dự án bao gồm:
\begin{itemize}
    \item Khảo sát hiện trạng, thu thập yêu cầu từ các sở/ngành và quận/huyện.
    \item Xây dựng phần mềm \gls{geoai}: module \gls{ai}, ứng dụng desktop.
    \item Kiểm thử chức năng từng module.
    \item Kiểm thử an toàn an ninh thông tin.
    \item Cài đặt, triển khai hệ thống.
    \item Đào tạo cán bộ quản lý và người dùng.
\end{itemize}

\subsection{Phân tích các yêu cầu của phần mềm}

\subsubsection{Yêu cầu chức năng của phần mềm}
 \begin{longtable}{|c|p{7.5cm}|c|c|}
\caption{Phân loại và tính toán số giao dịch của các yêu cầu chức năng} \label{tab:so_giao_dich} \\
\hline
\rowcolor{gray!25}
\textbf{TT} & \centering\textbf{Mô tả yêu cầu} & \centering\textbf{Phân loại (B/M/T)} & \textbf{Số giao dịch} \\ \hline
\endfirsthead

% Cấu hình phần đầu cho trang sau: CHỈ KẺ ĐƯỜNG NGANG, BỎ DÒNG TIÊU ĐỀ
\hline
\endhead

% Cấu hình phần chân bảng khi bị ngắt trang
\hline
\endfoot

% Cấu hình khi kết thúc bảng
\hline
\endlastfoot

1 & Quản lý người dùng và phân quyền & M & 6 \\ \hline
2 & Quản lý danh mục loại tài sản & M & 5 \\ \hline
3 & Quản lý danh mục khu vực địa lý & M & 5 \\ \hline
4 & Cấu hình hệ thống và bản đồ & B & 3 \\ \hline
5 & Xem nhật ký và kiểm toán & B & 3 \\ \hline
6 & Thêm tài sản mới trên bản đồ & M & 5 \\ \hline
7 & Chỉnh sửa thông tin tài sản & M & 4 \\ \hline
8 & Xóa tài sản & B & 3 \\ \hline
9 & Tìm kiếm và lọc tài sản trên bản đồ & M & 6 \\ \hline
10 & Import dữ liệu tài sản hàng loạt & M & 4 \\ \hline
11 & Xuất dữ liệu tài sản & M & 4 \\ \hline
12 & \gls{ai} phân tích ảnh và phân loại tài sản & T & 8 \\ \hline
13 & Tạo công việc bảo trì & M & 6 \\ \hline
14 & Cập nhật tiến độ bảo trì & M & 5 \\ \hline
15 & Theo dõi tiến độ bảo trì & B & 3 \\ \hline
16 & Tra cứu lịch sử bảo trì tài sản & B & 3 \\ \hline
17 & Xem bản đồ hạ tầng công khai & B & 3 \\ \hline
18 & Gửi phản ánh sự cố hạ tầng & M & 6 \\ \hline
19 & Tra cứu kết quả xử lý phản ánh & M & 4 \\ \hline
20 & Xem Dashboard tổng quan & B & 3 \\ \hline
21 & Xem báo cáo thống kê tài sản & B & 3 \\ \hline
22 & Xuất báo cáo và dữ liệu & B & 3 \\ 
\end{longtable} 
 \subsubsection{Yêu cầu Kiểm thử phần mềm}
\begin{itemize}
    \item Kiểm thử đơn vị (Unit Test) cho toàn bộ các function nghiệp vụ quan trọng, tỷ lệ coverage $\ge$ 80\%.
    \item Kiểm thử tích hợp giữa các module và với \gls{api} bên ngoài.
    \item Kiểm thử chức năng theo use case: mỗi use case nhập tối thiểu 10 bộ dữ liệu đầu vào.
    \item Kiểm thử hiệu năng: hệ thống đáp ứng 500 người dùng đồng thời, thời gian phản hồi $\le$ 3 giây.
    \item Kiểm thử bảo mật: quét lỗ hổng, kiểm thử xâm nhập.
    \item Kiểm thử trên bản đồ: xác thực tọa độ GPS, hiển thị đúng điểm trên bản đồ, render nhanh khi có 100.000+ điểm tài sản.
\end{itemize}

\subsubsection{Yêu cầu phi chức năng}

\textbf{a) Yêu cầu về cơ sở dữ liệu}
\begin{itemize}
    \item Sử dụng PostgreSQL với extension PostGIS để lưu trữ và xử lý dữ liệu không gian địa lý.
    \item Hỗ trợ truy vấn không gian địa lý: giao, tiếp xúc, bao phủ, khoảng cách.
    \item Backup tự động hàng ngày, lưu trữ 90 ngày, khôi phục trong vòng 2 giờ.
\end{itemize}

\textbf{b) Yêu cầu về bảo mật}
\begin{itemize}
    \item Xác thực đa yếu tố (MFA) cho tài khoản quản trị.
    \item Mã hóa dữ liệu truyền tải TLS 1.3, mã hóa dữ liệu lưu trữ AES-256.
    \item Phân quyền chi tiết theo vai trò và khu vực địa lý.
\end{itemize}

\textbf{c) Yêu cầu về giao diện người dùng}
\begin{itemize}
    \item Giao diện \gls{webgis} responsive, hỗ trợ đầy đủ trên máy tính, tablet và điện thoại.
    \item Bản đồ tương tác: zoom, pan, đo khoảng cách, chuyển đổi lớp bản đồ nền.
    \item Hỗ trợ tiếng Việt hoàn toàn, bao gồm tìm kiếm địa chỉ bằng tiếng Việt.
\end{itemize}

\textbf{d) Yêu cầu về tốc độ xử lý}
\begin{itemize}
    \item Tải bản đồ với 100.000+ điểm tài sản trong vòng 5 giây.
    \item \gls{ai} phân tích ngữ nghĩa và trả kết quả truy vấn không gian trong vòng 3 giây/lệnh.
    \item \gls{api} phản hồi trong vòng 2 giây cho 95\% các request thông thường.
\end{itemize}

\subsection{Danh mục các chuẩn, tiêu chuẩn áp dụng}
Hệ thống \gls{geoai} được thiết kế và triển khai tuân thủ nghiêm ngặt các tiêu chuẩn quốc gia và quốc tế:
\begin{itemize}
    \item \textbf{Tiêu chuẩn Kiến trúc:} Khung kiến trúc Chính phủ điện tử Việt Nam phiên bản 3.0.
    \item \textbf{Tiêu chuẩn Dữ liệu Không gian:} Tuân thủ các chuẩn của OGC (Open Geospatial Consortium) như WMS, WFS và GeoJSON cho việc trao đổi dữ liệu \gls{gis}.
    \item \textbf{Tiêu chuẩn An toàn thông tin:} Tuân thủ ISO/IEC 27001 và các tiêu chuẩn bảo mật ứng dụng web của OWASP (Open Web Application Security Project).
    \item \textbf{Tiêu chuẩn Mã hóa:} Sử dụng thuật toán AES-256 cho dữ liệu lưu trữ và chuẩn mã hóa TLS 1.3 cho dữ liệu truyền tải.
\end{itemize}

% --- CHUYỂN SANG PHẦN MỚI ---

\subsection{Đề xuất giải pháp kỹ thuật công nghệ}

\subsubsection{Kiến trúc hệ thống tổng thể}
\gls{geoai} áp dụng kiến trúc Monorepo (Turborepo) kết hợp Micro-services, phân chia rõ ràng các lớp: Frontend \gls{webgis} (Next.js), Backend \gls{api} (NestJS), Data Layer (PostgreSQL/PostGIS + Redis) và External \gls{ai} Services (Groq \gls{api}).

\subsubsection{Giải pháp \gls{gis} và Tìm kiếm dữ liệu}
\begin{itemize}
    \item Xử lý không gian: Sử dụng trực tiếp sức mạnh của PostgreSQL kết hợp PostGIS để thực hiện các truy vấn không gian phức tạp (ST\_DWithin, ST\_Distance) thông qua Prisma ORM, giúp tối ưu hiệu năng.
    \item Full-text Search \& Analytics: Tích hợp \textbf{Elasticsearch} làm công cụ tìm kiếm văn bản toàn văn (tên đường, tên tài sản) và phân tích log với tốc độ cao, giảm tải cho PostgreSQL trong các tác vụ tìm kiếm tự do.
    \item Bản đồ nền: Tích hợp OpenStreetMap thông qua thư viện Leaflet (React-Leaflet) trên frontend.
\end{itemize}

\subsubsection{Tích hợp Dữ liệu Bên thứ 3 (3rd Party Data)}
Hệ thống kết nối và tiêu thụ dữ liệu từ các nền tảng bên thứ ba để làm giàu hệ sinh thái thông tin:
\begin{itemize}
    \item \textbf{Overture Maps:} Tự động đồng bộ khối lượng lớn dữ liệu địa điểm, tiện ích công cộng (POI) và hình học (footprint) của các tòa nhà.
    \item \textbf{OpenStreetMap (OSM):} Khai thác mạng lưới đường giao thông (RoadSegment) để phục vụ cho các thuật toán đánh giá khả năng tiếp cận và chấm điểm tài sản.
\end{itemize}

\subsubsection{Giải pháp \gls{ai}/Machine Learning}
\begin{itemize}
    \item Spatial NLP Agent: Tích hợp mô hình ngôn ngữ lớn (\gls{llm}) Llama 3.3 thông qua Groq \gls{api} để dịch ngôn ngữ tự nhiên thành các truy vấn không gian phức tạp.
    \item Đánh giá Rủi ro Không gian (Risk Assessment): Tự động overlay dữ liệu tài sản với các lớp bản đồ rủi ro (ngập lụt, sạt lở) để cảnh báo.
    \item Asset Scoring: Thuật toán chấm điểm tài sản tự động dựa trên mức độ tiếp cận đường giao thông, tiện ích xung quanh (trường học, bệnh viện) và trừ điểm rủi ro.
\end{itemize}

\subsubsection{Giải pháp Backend \gls{api}}
\begin{itemize}
    \item Framework: NestJS (TypeScript/Node.js) - cung cấp kiến trúc module hóa chặt chẽ, dễ bảo trì và mở rộng.
    \item Database ORM: Prisma ORM cho an toàn kiểu dữ liệu (Type-safe) và quản lý schema.
    \item Cache: Redis qua nestjs/cache-manager để lưu trữ đệm các truy vấn không gian lặp lại và kết quả phân tích \gls{ai}.
\end{itemize}

\subsubsection{Giải pháp lưu trữ và bảo mật}
\begin{itemize}
    \item Object Storage: Sử dụng AWS S3 Protocol cho lưu trữ ảnh tài sản và báo cáo đính kèm.
    \item Bảo mật: Xác thực dựa trên JWT (JSON Web Token), mã hóa mật khẩu bằng bcrypt. Hệ thống phân quyền (RBAC) chi tiết đến từng module chức năng và lớp bản đồ.
\end{itemize}
\cleardoublepage